\section{Samples, averages, and finitary laws of large numbers}\label{sec:samples}

\subsection{Sample means}

Let \(X_0,\ldots,X_{n-1}\) be coordinates on a finite product law.  Their
sample mean is
\[
 \overline X_n=\frac1n\sum_{j<n}X_j
\]
for \(n>0\).

\begin{theorem}[Unbiased sample mean]\label{thm:sample-mean}
If each \(X_j\) has expectation \(\mu\), then
\[
 \E[\overline X_n]=\mu.
\]
No independence assumption is needed.
\end{theorem}

\begin{proof}
By finite linearity,
\[
 \E[\overline X_n]
 =\frac1n\sum_{j<n}\E[X_j]
 =\frac1n(n\mu)=\mu.
\]
In \(\ipra\), the identity \(\sum_{j<n}\mu=n\mu\) is proved by
quantifier-free induction on \(n\), and the theorem is kept in
cross-multiplied form \(n\E[\overline X_n]=n\mu\).
\end{proof}

\begin{theorem}[Variance of the independent sample mean]\label{thm:sample-var}
If the coordinates are independent and identically distributed with variance
\(\sigma^2\), then
\[
 \Var(\overline X_n)=\frac{\sigma^2}{n}.
\]
\end{theorem}

\begin{proof}
By independent variance additivity,
\(\Var(\sum_{j<n}X_j)=n\sigma^2\).  Homogeneity of variance,
\(\Var(aX)=a^2\Var(X)\), follows from its definition.  Therefore
\[
 \Var(\overline X_n)=\frac1{n^2}n\sigma^2=\frac{\sigma^2}{n}.
\]
All calculations are rational when the underlying finite law is rational.
\end{proof}

\begin{theorem}[Expected residual sum of squares]\label{thm:residual}
For independent identically distributed coordinates with variance
\(\sigma^2\),
\[
 \E\left[\sum_{j<n}(X_j-\overline X_n)^2\right]=(n-1)\sigma^2.
\]
Consequently, for \(n>1\),
\(\frac1{n-1}\sum_{j<n}(X_j-\overline X_n)^2\) is an unbiased estimator of
\(\sigma^2\).
\end{theorem}

\begin{proof}
The pointwise polynomial identity
\[
 \sum_{j<n}(X_j-\overline X_n)^2
 =\sum_{j<n}(X_j-\mu)^2-n(\overline X_n-\mu)^2
\]
is proved by expanding both sides and using
\(\sum_j(X_j-\overline X_n)=0\).  Take expectations.  The first term has
expectation \(n\sigma^2\); the second has expectation
\(n\Var(\overline X_n)=\sigma^2\).  Subtract.
\end{proof}

\subsection{Weak law as a finite quantitative theorem}

\begin{theorem}[Chebyshev weak-law bound]\label{thm:weak-law}
Under the assumptions of \Cref{thm:sample-var}, for rational
\(\varepsilon>0\),
\[
 \Pp(\abs{\overline X_n-\mu}\ge\varepsilon)
 \le\frac{\sigma^2}{n\varepsilon^2}.
\]
Given rational \(\delta>0\), any natural \(n\) with
\(n\varepsilon^2\delta\ge\sigma^2\) guarantees probability at most
\(\delta\).
\end{theorem}

\begin{proof}
Apply Chebyshev to \(\overline X_n\), insert
\(\Var(\overline X_n)=\sigma^2/n\), and clear the positive denominators.  The
last sentence is merely transitivity of rational order.  It is a direct
\(\ipra\) contract: the input supplies \(n,\varepsilon,\delta,\sigma^2\) and
a checker verifies one natural inequality.
\end{proof}

\subsection{Kolmogorov's finite maximal inequality}

For a finite independent centered sequence put
\(S_k=\sum_{j<k}X_j\).  Every product sample space is finite, so first-crossing
events are decidable by bounded search.

\begin{theorem}[Finite Kolmogorov maximal inequality]\label{thm:kolmogorov}
Let \(X_0,\ldots,X_{n-1}\) be independent finite rational observables with
\(\E[X_j]=0\).  For rational \(a>0\),
\[
 \Pp\left(\max_{1\le k\le n}\abs{S_k}\ge a\right)
 \le\frac{\sum_{j<n}\Var(X_j)}{a^2}.
\]
\end{theorem}

\begin{proof}
Let \(A_k\) be the event that \(k\) is the first index with
\(\abs{S_k}\ge a\).  The \(A_k\) are disjoint and their union is the maximal
event.  On \(A_k\), \(S_k^2\ge a^2\), so
\[
 a^2\Pp(A_k)\le\E[S_k^2\ind_{A_k}].
\]
Write \(S_n=S_k+(S_n-S_k)\).  Since \(A_k\) depends only on the first \(k\)
coordinates and the tail increment is independent of them with mean zero,
\[
 \E[S_k(S_n-S_k)\ind_{A_k}]=0.
\]
Also \(\E[(S_n-S_k)^2\ind_{A_k}]\ge0\).  Hence
\[
 \E[S_n^2\ind_{A_k}]\ge\E[S_k^2\ind_{A_k}]
 \ge a^2\Pp(A_k).
\]
Sum over the disjoint first-crossing events:
\[
 a^2\Pp\left(\max_{k\le n}\abs{S_k}\ge a\right)
 \le\E\left[S_n^2\sum_k\ind_{A_k}\right]
 \le\E[S_n^2].
\]
The final second moment is \(\sum_{j<n}\Var(X_j)\) by independence and
centering.  Every event and every sum in this proof is finite.  The only
induction is the construction of partial sums and the bounded search for the
first crossing.
\end{proof}

\subsection{A finitary strong-law contract}

A qualitative statement ``\(\overline X_n\to\mu\) almost surely'' quantifies
over an infinite sequence and does not by itself expose a usable rate.  The
following theorem is the precise finite-error content needed by many
applications.

\begin{theorem}[Dyadic eventual-uniform bound]\label{thm:dyadic-slln}
Let \(X_0,X_1,\ldots\) be presented as an iid finite rational law with mean
\(\mu\) and variance \(\sigma^2\).  For naturals \(K\) and \(H\ge2^K\), and
rational \(\varepsilon>0\),
\[
 \Pp\left(\exists n\in[2^K,H]\;:
   \abs{\overline X_n-\mu}\ge 2\varepsilon\right)
 \le \frac{4\sigma^2}{\varepsilon^2 2^K}.
\]
Consequently, if
\(4\sigma^2\le\delta\varepsilon^2 2^K\), then with probability at least
\(1-\delta\), every \(n\in[2^K,H]\) satisfies
\(\abs{\overline X_n-\mu}<2\varepsilon\).
\end{theorem}

\begin{proof}
Center the variables and write \(S_n=\sum_{j<n}(X_j-\mu)\).  Partition
\([2^K,H]\) into dyadic blocks \([2^r,2^{r+1})\).  If for some
\(n\in[2^r,2^{r+1})\) we have \(\abs{S_n}\ge2\varepsilon n\), then certainly
\(\max_{k\le2^{r+1}}\abs{S_k}\ge2\varepsilon2^r\).  Kolmogorov's inequality
gives
\[
 \Pp(\text{bad block }r)
 \le\frac{2^{r+1}\sigma^2}{(2\varepsilon2^r)^2}
 =\frac{\sigma^2}{2\varepsilon^2 2^r}.
\]
The union bound over \(r\ge K\) up to the last block is at most the infinite
geometric majorant
\[
 \sum_{r=K}^{\infty}\frac{\sigma^2}{2\varepsilon^2 2^r}
 =\frac{\sigma^2}{\varepsilon^2 2^K},
\]
which is stronger than the displayed constant \(4\).  We retain \(4\) to
absorb endpoint conventions, non-power-of-two horizons, and a version obtained
by controlling block increments rather than maxima from zero.  Formally, for a
fixed finite \(H\), only a finite geometric sum is used, and it is bounded by
\(2^{1-K}\) through an elementary induction.

Finally, \(\abs{\overline X_n-\mu}=\abs{S_n}/n\).  The complement of the bad
event therefore gives the simultaneous bound.  The theorem is a family of
finite \(\ipra\) statements uniform in \(K,H,\varepsilon,\delta\).
\end{proof}

\begin{remark}[Relation to the classical strong law]
The family in \Cref{thm:dyadic-slln} has a tail bound independent of the finite
horizon \(H\).  Metamathematically, letting \(H\) grow and then \(K\) grow
recovers the usual almost-sure convergence interpretation.  Inside \(\ipra\),
the directly usable theorem is the finite contract: given an error tolerance,
failure budget, and finite observation horizon, compute and verify a burn-in
index.  This avoids silently importing countable additivity or an infinitary
``eventually'' predicate.
\end{remark}
